\subsection{Surface relief: procedural bump}
\code{surfaceHeight(q)} in \code{ShaderSources.cpp} gives each procedural material a height field in the same triplanar coordinates as its albedo (concrete pores 1.5\,mm, bark ridges, rock, planks, corrugated ribs 9\,mm, HESCO mesh and bags, asphalt
aggregate, camo paint chips, steel seams $-3$\,mm and rivets $+2$\,mm). Listing~\ref{lst:Bump} shows \code{surfaceBump(P,N)}, which takes central differences $\partial h/\partial u,\partial h/\partial v$ ($\epsilon=4$\,mm) and tilts the normal,
$\vect N_b=\mathrm{normalize}\bigl(\vect N-(h_u\vect T+h_v\vect B)\,k\bigr)$, where $\vect T,\vect B$ is the tangent frame of the dominant face and $k$ fades out from 12 to 40\,m (no far shimmer). Because the Gouraud and Flat paths evaluate lighting before the fragment
stage, the light terms are \emph{re-aimed} by the Lambert ratio
\begin{equation}
k_L=\mathrm{clamp}\!\left(\frac{\max(\vect N_b\!\cdot\!\vect L,0)+0.03}{\max(\vect N\!\cdot\!\vect L,0)+0.03},\,0.3,\,2\right),
\end{equation}
applied to the moon and searchlight diffuse terms ($k_L$) and specular terms ($k_L^2$); a cavity factor $\mathrm{clamp}(0.6+0.4\,\vect N_b\!\cdot\!\vect N+0.5(N_{b,y}-N_y),0.6,1.15)$ darkens pits that face away from the sky, and the bumped normal also feeds the SSR surface buffer. The bump is active
only in HDR on HIGH and above (\code{uGroundDetail}); screens, foliage and the ground have their own code. \lstBump
\begin{figure}[h]
\centering
\piclab{bump_on_off}{\linewidth}
\caption{Bump off (left) and on (right) for corrugated metal, HESCO and steel plate: ribs, mesh and rivets catch the light (in-program render).}
\end{figure}
The planned \emph{reflection probe} (a low-resolution sky cubemap for metal and glass) is \textbf{not} implemented; reflections still come from SSR and the analytic sky.
